% problem-000097 3 双氧物种电子与晶体
\section*{3. 双氧物种电子与晶体}
（图：）

图中⿊圈部分为：

\subsection*{3-1}
MOF-5的晶体结构如上图所⽰，写出MOF-5的化学式。

\subsection*{3-2}
MOF-5具有较⼤的孔洞，制备过程中溶剂分⼦会占据孔洞，已知MOF-X的密度为 $1 . 1 5 3 \mathrm { g c m ^ { - 3 } } _ { \mathrm { c } }$

MOF-X的晶胞参数如下： $a = 2 5 6 6 . 9 0 \mathrm { p m }$ $V = 1 . 6 9 1 3 2 \times 1 0 ^ { 1 0 } \mathrm { p m } ^ { 3 }$ ，� =8；含氯量为2.42%。推断⼀个正当晶胞中溶剂分⼦的数⽬与种类。

\subsection*{3-3}
MOF-5具有可以容纳直径为1510pm球状分⼦的孔洞，孔洞的体积约占晶体总体积的55%；通过加热等⼿段完全除去溶剂分⼦的MOF-5能稳定存在。按照堆积理论这样的结构很不稳定，但实验发现如此⾼孔隙率的MOF-5却能稳定存在，说明原因。
